\section{The full Brown--Henneaux metric quotient}\label{sec:brown-henneaux-theorem}

\subsection{Reducibility and the endpoint-fixing section}\label{subsec:psl-section}

Consider one chiral Brown--Henneaux generator labelled by $f\in H^\sigma(S^1)$. The metric perturbation is unchanged when $f$ is shifted by the three global conformal modes $1$, $\cos\phi$, and $\sin\phi$, because their bulk lifts are exact AdS Killing fields. Endpoint values of $f$ are therefore not functions on the metric tangent space.

For a symmetric interval $A=(-a,a)$ with $0<a<\pi/2$, define

\begin{align}
q_\pm[f]=f(\pm a)
\end{align}

and the bounded section

\begin{align}
P_{\rm PSL}f
=f-\frac{q_++q_-}{2\cos a}\cos\phi
-\frac{q_+-q_-}{2\sin a}\sin\phi . \label{eq:psl-section}
\end{align}

Then

\begin{align}
(P_{\rm PSL}f)(\pm a)=0,
\qquad
P_{\rm PSL}^2=P_{\rm PSL},
\qquad
h[P_{\rm PSL}f]=h[f]. \label{eq:psl-properties}
\end{align}

The residual global direction that already fixes both endpoints is the interval modular weight

\begin{align}
w_A(\phi)=\frac{\cos\phi-\cos a}{\sin a}.
\end{align}

Consequently,

\begin{align}
H^\sigma(S^1)/\mathfrak{sl}(2,\mathbb R)
\simeq
\ker(q_+,q_-)/\operatorname{span}\{w_A\}. \label{eq:quotient-isomorphism}
\end{align}

Equation (\ref{eq:psl-section}) is not the spectral projector that adds charged $m\geq2$ modes to create a special fixed-anchor subspace. It changes only the global-Killing representative and covers the full metric quotient. For other interval charts, one chooses a nonsingular pair of global conformal complements instead of dividing by a vanishing $\cos a$ or $\sin a$.

\subsection{The proper Hollands--Wald section}\label{subsec:hw-section}

On the complete RT geodesic, the Jacobi inverse of $J=-D_u^2+1$ is

\begin{align}
G_J(u,u')=\frac12e^{-|u-u'|}.
\end{align}

For endpoint-zero data, (\ref{eq:jacobi}) therefore has a unique decaying normal solution. For a general representative, the growing part of $-\zeta_\perp$ is cancelled by a Jacobi-homogeneous correction. In a parallel two-component normal frame it is

\begin{align}
H^{(0)}&=-\frac{q_+e^u+q_-e^{-u}}{\sqrt2},\\
H^{(1)}&=\frac{-q_+e^u+q_-e^{-u}}{\sqrt2}. \label{eq:homogeneous-hw}
\end{align}

The correction (\ref{eq:homogeneous-hw}) is precisely the normal restriction of the global Killing field removed by $P_{\rm PSL}$. The resulting vector has the proper falloffs

\begin{align}
(v^t,v^r,v^\phi)
=(O(r^{-2}),O(r^{-1}),O(r^{-2}))
\end{align}

at both anchors. Extending its normal two-jet through a fixed-width Fermi collar defines a bounded weak Sobolev section $P_{\rm HW}$ on the regularity class used below. At every finite regulator,

\begin{align}
\Delta_Xg(P_{\rm HW}h,-V[h])=h. \label{eq:hw-pullback}
\end{align}

Thus every bulk, GHY, counterterm, Harlow--Wu, Hayward, anchor, and embedding term in the complete pulled-back action equals its fixed-section value. Equation (\ref{eq:hw-pullback}) is the action-level reason that the HW collar introduces no new flux estimate.

\subsection{Linked regulators and slice covariance}\label{subsec:slice-covariance}

Let $M$ be a frequency cutoff and parameterize the regulators by

\begin{align}
L=M^\alpha,
\qquad
y:=\tanh\frac\epsilon2=\frac12M^{-\beta},
\qquad
R=M^\gamma . \label{eq:linked-schedule}
\end{align}

The complete list of exponents is derived in Appendix \ref{app:regulator-estimates}. Every displayed monomial tends to zero in the open region

\begin{align}
\alpha>\frac72,
\qquad
\beta>0,
\qquad
\gamma>2\alpha+3. \label{eq:formal-exponent-region}
\end{align}

The $L^4/R^2$ Cauchy-corner term is the reason for the last, stronger inequality. The exact complete finite-wall wall kernel has presently been proved uniformly on $y=1/(2L)$, not for arbitrary unrelated $y(L)$. The proved regulator class used in the theorem is therefore

\begin{align}
\alpha>\frac72,
\qquad
\beta=\alpha,
\qquad
\gamma>2\alpha+3. \label{eq:proved-exponent-region}
\end{align}

The earlier numerical choice $(\alpha,\beta,\gamma)=(10,10,30)$ is one member. It gives $O(M^{-14})$ for the transition term, $O(M^{-13})$ for the complete compensated finite wall, $O(M^{-34}+M^{-14})$ for the compensated Cauchy corner, and $O(M^{-34})$ for the compensated outer joint. The outer Brown--York source curvature vanishes exactly.

This also proves schedule independence within (\ref{eq:proved-exponent-region}): the difference between two regulated forms is bounded by the sum of their two vanishing error majorants. The same triangle argument proves taper independence for any two taper families which separately satisfy the uniform transition, wall, and outer-corner estimates of Appendix \ref{app:regulator-estimates} and approach the same local kernel. It is not a claim that every smooth taper satisfies those hypotheses. Extending the complete wall majorant to an independent $\beta$ would promote (\ref{eq:formal-exponent-region}) from a monomial compatibility region to a proved two-parameter theorem.

The construction is covariant under boundary time translation. For left movers, let $U_tf(\phi)=f(\phi-t)$. The endpoint, PSL, and HW sections at time $t$ are obtained by conjugation:

\begin{align}
P_{{\rm PSL},t}=U_{-t}P_{\rm PSL}U_t,
\qquad
P_{{\rm HW},t}=U_{-t}P_{\rm HW}U_t. \label{eq:time-sections}
\end{align}

The Brown--Henneaux phases have unit modulus, so the estimates are uniform in $t$. Vanishing total flux makes time translation a symplectomorphism between the transported fibres. Time translation itself is a physical AdS isometry; the fixed/HW section change inside one fibre is proper gauge.

\subsection{Main theorem}\label{subsec:main-theorem}

\textbf{Theorem (classical finite-action subregion identity).} Let $(M,G)$ be vacuum AdS$_3$, let $A=(-a,a)$ be a boundary interval whose RT geodesic $\gamma_A$ is the bifurcation surface of the exact modular Killing field $\xi_A$, and impose Brown--Henneaux source boundary conditions. Use the finite-action CPS convention of Section \ref{sec:finite-action-cps}, the moving pullback of Section \ref{sec:moving-boundary}, and the proved linked regulator class (\ref{eq:proved-exponent-region}). Then, for each chirality and every

\begin{align}
[f]\in H^\sigma(S^1)/\mathfrak{sl}(2,\mathbb R),
\qquad \sigma>\frac52,
\end{align}

the second-order charge satisfies

\begin{align}
\delta^2H_{\xi_A}^\infty
&=\frac{s_{\xi_A}}{2\pi}
\delta^2\!\left(\frac{A[g,\gamma_A]}{4G}\right)
+E_{{\rm can,p}}(P_{\rm HW}h) \label{eq:main-theorem-hw}\\
&=\frac{s_{\xi_A}}{2\pi}
\delta^2\!\left(\frac{A[g,\gamma_A]}{4G}\right)
+E_{{\rm can,p}}[h]
+\int_{\gamma_A}\Upsilon_{\rm p}[h,V[h]]. \label{eq:main-theorem-separated}
\end{align}

Both expressions are independent of the global-Killing representative and of the transported Cauchy slice. They are also independent of an admissible ambient extension map for $V[h]$ when two extensions have the same prescribed cut data and asymptotic support and differ by a vector in $\mathfrak g_{\rm proper}$. A source-changing or charged extension is not covered by this statement. The real Brown--Henneaux space is the orthogonal direct sum of the two chiral quotients, and the two quadratic forms add.

The regularity $\sigma>5/2$ is required for the separated point-anchor chart in (\ref{eq:main-theorem-separated}), where second-derivative endpoint traces occur. It is not the optimal regularity of the combined quadratic form. The theorem does not quotient the charged Virasoro modes $m\geq2$, introduce a negative-Virasoro frame, or assume an auxiliary edge cotangent sector.
